\clearpage\begingroup\SourceLeaf{611}
\ArithmeticHead{56}{L}{PHILIPP MAENNCHEN, GAUSS ALS ZAHLENRECHNER.}
\renewcommand{\thefootnote}{\arabic{footnote})}\setcounter{footnote}{0}
\noindent\hspace*{\parindent}Nimmt also $c$ um 1 ab, so nimmt $e$ um 4 ab. Das gleiche zeigt sich
bei der weiteren Behandlung der Gruppe:
\[
\begin{gathered}
n_{J+152}\equiv n_{J+40}\equiv n_{J+12}\equiv n_J+12+3\equiv n_J+1\; (\mathrm{mod}\ 7)\\
e_{J+152}\equiv e_J-1\equiv e_J-8\; (\mathrm{mod}\ 7)\quad\text{usw.}
\end{gathered}
\]
\noindent Ergebnis: $e$ ist eine lineare Funktion von $c$, und zwar ist
\[
e\equiv4c+\sigma\; (\mathrm{mod}\ 7),
\]
\noindent wobei $\sigma$ von $c$ unabhängig ist.

Fassen wir die drei Ergebnisse zusammen, so folgt:
\par\noindent $e$ ist eine lineare Funktion von $b$, $c$, $d$, und zwar ist
\[
e\equiv2b+4c+6d+\tau\; (\mathrm{mod}\ 7),
\]
\noindent wo $\tau$ eine Konstante ist. Diese Konstante kann man nun leicht bestimmen,
indem man entweder für irgend ein Jahr den Wochentag des Frühlingsvollmonds berechnet, oder indem man in einer Ostertabelle für irgend ein Jahr
den Ostertermin abliest und in die obigen Formeln einsetzt. Für das Jahr
980 z. B. hat man $b=0$, $c=0$, $d=14$, $O=42$, also $e=6$, somit
\[
6\equiv0+0+84+\tau\; (\mathrm{mod}\ 7).
\]
\noindent Also ist $\tau\equiv6$, und die Formel lautet endgültig:
\[
e\equiv2b+4c+6d+6\; (\mathrm{mod}\ 7),
\]

Damit ist eine Berechnungsmethode für das julianische Osterfest gewonnen, genau in der Form, wie sie \textsc{Gauss} geprägt hat. Sie hat der \textsc{Lambert}schen Formel gegenüber den bedeutsamen Vorteil, dass sie geeignet ist,
sich sowohl den veränderten Verhältnissen des gregorianischen Kalenders, als
auch sogar den sonderbaren Ausnahmebestimmungen für das gregorianische
Osterfest anzupassen.

Bei dieser Anpassung werden nur die Formeln für $d$ und $e$ etwas abgeändert, alles übrige bleibt. Wir erhalten nach \textsc{Gauss}:
\[
d\equiv19a+M\; (\mathrm{mod}\ 30)
\]
\noindent und
\[
e\equiv2b+4c+6d+N\; (\mathrm{mod}\ 7).
\]
\noindent Die Werte $M$ und $N$ sind im Laufe eines Jahrhunderts konstant; sie sind
nur beim Übergang von einem Jahrhundert in das andere gewissen Veränderungen unterworfen. Die Veränderung von $N$ ist bedingt durch die bei der
\clearpage\SourceLeaf{612}\ArithmeticHead{57}{R}{X. CHRONOLOGISCHE ARBEITEN.}
\noindent gregorianischen Reform neu eingeführte sogenannte \textls[120]{Sonnengleichung der
Epakte}. Diese bewirkt, dass infolge des Ausfalls des Schalttags in den Jahren
1700, 1800, 1900, 2100, .\ .\ .\ . die Epakte um 1 kleiner wird (Epaktensprung),
und somit das Datum des Frühlingsvollmonds um einen Tag vorrückt. Die
Veränderung von $M$ ist ebenfalls durch diese Sonnengleichung beeinflusst, aber
ausserdem auch durch die \textls[120]{Mondgleichung}. Diese Grösse entsteht dadurch,
dass der 19jährige Metonische Zyklus nicht genau ein ganzzahliges Vielfaches
des synodischen Monats beträgt, sondern dass nach Ablauf eines solchen
Zyklus die Epakte um $0{,}0609$ Tage zu klein geworden ist, d. h. dass der wirkliche Neumond nun um $0{,}0609$ Tage früher eintritt, als der berechnete \textls[120]{ekklesiastische} Neumond. Nach Verlauf von $16{,}42$ solcher Zyklen würde
dieser Fehler zu einem vollen Tage angewachsen sein, und man hätte daher
nach $16{,}42\mathord{.}19=312$ Jahren die Epakte um 1 erhöhen müssen, um wieder
Übereinstimmung zu erhalten. Dies geschieht erst seit der gregorianischen
Reform, und zwar derart, dass man alle 300 Jahre die Epakte um 1 erhöht.
Man wählt diejenigen Jahrhundertanfänge, deren Hunderter durch 3 teilbar
sind, also 1800, 2100 usw. Diese Erhöhung der Epakte infolge der Mondgleichung bedingt einen Rückgang der Ostergrenze um einen Tag, während
der Rückgang der Epakte infolge der Sonnengleichung eine Erhöhung der
Ostergrenze um einen Tag zur Folge hat. Das sind die chronologischen
Daten, die der Bestimmung von $M$ und $N$ zugrunde liegen.

Wir wollen mit \textsc{Gauss} die Anzahl der Hunderter eines gegebenen Jahrhunderts mit $k$ bezeichnen und die in der Osterformel auftretende Zahl $M$ aus den Summanden $k$ und $\nu$ zusammensetzen. Ferner wollen
wir die Werte von $k$ von 15 bis 24 untereinanderschreiben und neben jedes $k$ den zugehörigen Wert von $M$ setzen unter Berücksichtigung von Sonnen- und Mondgleichung. In die
folgende Spalte schreiben wir die zugehörigen
Werte von $\nu$, und neben diese noch die Werte
$p$ und $q$, welche die ganzzahligen Bestandteile
von $\frac{k}{4}$ und $\frac{k}{3}$ darstellen. Denn diese ganzzahligen Bestandteile erfahren alle 4 bezw. 3 Jahre eine Veränderung; der
\begin{center}\setlength{\tabcolsep}{5pt}\renewcommand{\arraystretch}{1.15}
\begin{tabular}{c|l|c|c|c}\hline\hline
$k$&\multicolumn{1}{c|}{$M$}&$\nu$&$p$&$q$\\\hline\hline
15&$22\hphantom{{}+0-0}=22$&7&3&5\\
16&$22+0-0=22$&6&4&5\\
17&$22+1-0=23$&6&4&5\\
18&$23+1-1=23$&5&4&6\\
19&$23+1-0=24$&5&4&6\\
20&$24+0-0=24$&4&5&6\\
21&$24+1-1=24$&3&5&7\\
22&$24+1-0=25$&3&5&7\\
23&$25+1-0=26$&3&5&7\\
24&$26+0-1=25$&1&6&8
\end{tabular}\end{center}
\par\noindent X 2 Abh. 6.\hfill 8
\clearpage\SourceLeaf{613}\ArithmeticHead{58}{L}{PHILIPP MAENNCHEN, GAUSS ALS ZAHLENRECHNER.}
\setcounter{footnote}{0}
\noindent Einfluss der Sonnengleichung bedingt alle 4 Jahre \textls[120]{keine} Veränderung, der
der Mondgleichung alle 3 Jahre eine Veränderung. Man erkennt auf den
ersten Blick, dass $\nu+p+q=15$ ist, und somit die Richtigkeit der \textsc{Gauss}schen Formel
\[
M=k+15-p-q.
\]
\noindent Der Wert $M=22$ für $k=15$ ergibt sich leicht, wenn man für $k=15$ eine
Jahreszahl sucht, für die der früheste Ostertermin (22. März) eintrat. Das
war 1598. In diesem Fall musste $d=e=0$ sein, und so ergibt sich leicht
$M=22$ und $N=2$.

Für $N$ erhalten wir folgende Tabelle
\begin{center}\setlength{\tabcolsep}{5pt}\renewcommand{\arraystretch}{1.15}
\begin{tabular}{c|l|c}\hline\hline
$k$&\multicolumn{1}{c|}{$N$}&$p$\\\hline\hline
15&2&3\\
16&$2+0=2$&4\\
17&$2+1=3$&4\\
18&$3+1=4$&4\\
19&$4+1=5$&4\\
20&$5+0=5$&5\\
21&$5+1=6$&5\\
22&$6+1=7$&5\\
23&$7+1=8$&5\\
24&$8+0=8$&6
\end{tabular}\end{center}
\noindent Hier sieht man sofort:
\[
\begin{gathered}
k-N-p=10\\
N=k-p-10,
\end{gathered}
\]
\noindent oder, da wir nur den Rest von $N$ $(\mathrm{mod}\ 7)$ brauchen:
\[
N\equiv k-p+4\; (\mathrm{mod}\ 7).
\]

Damit ist nun die allgemeine Regel entwickelt, und
es wäre höchstens noch dazu zu bemerken, dass \textsc{Gauss}
sich dabei zunächst an eine genauere Festsetzung nicht
gekehrt hat. Da nämlich die Mondgleichung alle 300
Jahre berücksichtigt wird statt alle 312 Jahre, so würde
schon nach 2400 Jahren die Korrektur 8 mal angebracht sein, während diese
Korrektur nahezu für 2500 Jahre ausreicht. Von \textsc{Français} (1813) und von
\textsc{Tittel} (1816) darauf aufmerksam gemacht, berichtigte er die ursprüngliche
Fassung in der folgenden Weise:

»$p$ wird bestimmt als Quotient bei der Division von $8k+13$ mit 25«\footnote{Werke VI, S. 73, am Schluss.}.
\par\noindent Warum hat wohl \textsc{Gauss} nicht von vornherein die genauere Bestimmung von
$p$ gewählt? Es liegen zwei Möglichkeiten vor, dies zu erklären. Die eine
Möglichkeit ist die, dass er bei der Aufstellung seiner Formel sich gar nicht,
oder möglichst wenig, um die Regeln der Epaktenrechnung kümmerte, sondern dass er auch hier die Formel aus den Tabellen herauslas. Die andere,
wahrscheinlichere, ist wohl die, dass ihm die fernliegenden Jahrhunderte gleich-
\clearpage\SourceLeaf{614}\ArithmeticHead{59}{R}{X. CHRONOLOGISCHE ARBEITEN.}
\setcounter{footnote}{0}
\noindent gültig waren, da er überzeugt war, dass auch die gregorianische Reform nicht
für alle Zeiten bleiben würde. Diesem Gedanken gibt er auch Ausdruck am
Schluss des Aufsatzes: \textit{Noch etwas über die Bestimmung des Osterfestes} (Werke
VI, S. 82 ff., Braunschw. Magazin 1807, Sept. 12).

Jetzt muss noch den beiden eigentümlichen Ausnahmebestimmungen Rechnung getragen werden, und es mag für \textsc{Gauss} besonders reizvoll gewesen sein,
die beiden willkürlichen Festsetzungen über die Lage des Osterfestes zahlentheoretisch zu erfassen.

Die eine Bestimmung kann kurz so ausgesprochen werden: Der 25. April
ist der äusserste Ostertermin. Diese Bestimmung ist für das julianische Osterfest automatisch erfüllt, da der 49. März die äusserste Ostergrenze ist. Bei
der Formel für den gregorianischen Kalender ist sie dagegen nicht automatisch
erfüllt, da ja $M$ nicht mehr konstant ist. Man braucht aber nur folgende
Fassung zu wählen, die \textsc{Gauss}\footnote{Werke VI, S. 79.} vorgeschlagen hat: »Ergibt sich
nach der Rechnung der 26. April als Ostertermin, so fällt Ostern
tatsächlich auf den 19. April«, ein Fall, der im Jahre 1609 eintrat und 1981\footnote{Werke VI, S. 79, Zeile 11 heisst es statt 1981 irrtümlich 1989; die richtige Zahl findet sich z. B.
ebenda S. 85, Zeile 15 v. u.} wieder stattfinden wird.

Komplizierter ist die zweite willkürliche Bestimmung: Der
Frühlingsvollmond darf in einem Metonischen Zyklus nicht
zweimal auf denselben Tag fallen. Vielleicht geschah diese
Bestimmung auch in Anlehnung an die Tatsache, dass im
julianischen Kalender diese Eigenschaft automatisch erfüllt
war, während in den meisten Zyklen Doubletten von Osterterminen zu finden sind. Unter Metonischen Zyklen im gregorianischen Kalender muss man freilich nur solche verstehen, die
keinen Übergang in ein neues Jahrhundert enthalten, bei dem $\nu$
von Null verschieden ist; andernfalls wären nämlich Doubletten
nicht zu vermeiden. Wir wollen einen solchen Zyklus in allgemeiner Form hinschreiben. Zu jedem $a$ setzen wir den entsprechenden Wert von $d$, und da $d+21$ die Ostergrenze darstellt, so entspricht einer Doublette für die Ostergrenze auch
eine solche für $d$.
\begin{center}\setlength{\tabcolsep}{7pt}\renewcommand{\arraystretch}{1.12}
\begin{tabular}{c|l}\hline\hline
$a$&\multicolumn{1}{c}{$d$}\\\hline\hline
0&$d_0$\\
1&$d_0+19$\\
2&$d_0+8$\\
3&$d_0+27$\\
4&$d_0+16$\\
5&$d_0+5$\\
6&$d_0+24$\\
7&$d_0+13$\\
8&$d_0+2$\\
9&$d_0+21$\\
10&$d_0+10$\\
11&$d_0+29$\\
12&$d_0+18$\\
13&$d_0+7$\\
14&$d_0+26$\\
15&$d_0+15$\\
16&$d_0+4$\\
17&$d_0+23$\\
18&$d_0+12$
\end{tabular}\end{center}
\par\hfill $8^{\ast}$
\clearpage\endgroup
